\begin{frame}{비용함수: 평균제곱오차}
\begin{center}\begin{varwidth}{0.92\textwidth}\usebeamercolor[fg]{normal text}
  
  \[
  J(\boldsymbol{w}) = \frac{1}{2m}\sum_{i=1}^m \left(h_{\boldsymbol{w}}(\boldsymbol{x}^{(i)}) - y^{(i)}\right)^2
  \]
  
  \begin{itemize}
    
    
    \item \textbf{제곱은 두 문제를 동시에 해결}: 항상 양수 + 매끄럽게
      미분 가능
    \item 앞의 $\frac12$ 은 미분할 때 지수 2와 약분되라고 붙인 것.
      \textbf{평가 지표}로 쓸 때는 $\frac12$ 없이
      $\text{MSE} = \frac{1}{m}\sum_i (\hat y^{(i)} - y^{(i)})^2$,
      \ $\text{RMSE} = \sqrt{\text{MSE}}$ (원래 $y$ 와 같은 단위)
  \end{itemize}
\end{varwidth}\end{center}
\note{\begin{itemize}
\item $m$ 개의 학습 데이터에 대해, 모델이 얼마나 틀렸는가:
\item \textbf{왜 오차를 제곱해서 더하는가?}
\item 오차를 그냥 더하면 양수·음수 오차가 \textbf{상쇄} ---
      엉망인 모델도 ``평균 오차 0'' 가능
\item 절댓값은 상쇄는 막지만 미분이 \textbf{불연속} $\to$ 최적화가 까다로움
\end{itemize}}
\end{frame}
